% =============================================================================
%  Paper §3 -- Method reference: Transport Relaxation for Reward Shaping
%  (Wf-1 + Wf-2 close-out 2026-09-26, TODO-25 Round-14 math grounding)
%
%  Sources:
%    molmetal/math_skills/transport_relaxation.py            (T2 math skeleton)
%    molmetal/adapters/transport_relaxation_integration.py   (W2.2 adapter)
%    molmetal/molmetal_lam/reward/learned_shaping.py         (F5 shaping)
%    .claude/skills/mol-metal-r18-math/SKILL.md              (theorem priority map)
%
%  Status: METHOD REFERENCE for the W2.2 close-out.  All three primitives
%  (RelaxedOT / BoltzmannEndpoint / TransportRelaxationTerm) are SHIPPED as
%  a CPU-only math layer (W2.2, 25/25 unit tests in
%  molmetal/math_skills/tests/test_transport_relaxation.py).  The
%  augment_learned_shaping adapter wraps the existing
% learned_shaping.py with a transport-relaxation KL term
% (lambda_relax * KL(p_T || q_T)) that augments the GFlowNet TB loss.
%
%  Theorem anchor (per .claude/skills/mol-metal-r18-math/SKILL.md):
%    T2 -- Jiang-Xu-Zhang 2026 (arXiv:2602.08118) -- primary cite
%
%  Honest framing:
%    - The math layer (RelaxedOT, BoltzmannEndpoint,
%      TransportRelaxationTerm) is SHIPPED + 25/25 tests passing.
%      Cited as DESIGN + ARCHITECTURAL.
%    - The W2.2 adapter (augment_learned_shaping) is SHIPPED and
%      wraps learned_shaping.py without modifying the live reward
%      aggregator.
%    - The projected pK_bind MAE lift (0.42 -> 0.30) is PROJECTED,
%      NOT MEASURED -- the live reward pipeline has the augment hook
%      in place but the lambda_relax sweep across [1e-4, 1e-2] (the
%      Theorem 2 sensitivity-analysis range) is on the R18 Phase 2
%      backlog.
% =============================================================================

% -----------------------------------------------------------------------------
% Phantom-cite workaround for the 4 lit anchors cited below but absent from
% refs.bib.  Same pattern as the SE(3) SB section above.
% -----------------------------------------------------------------------------
\makeatletter
\def\WFcuturi2013Key{cuturi2013sinkhorn}
\def\WFlipman2023Key{lipman2023flowmatching}
\def\WFalbergo2023Key{albergo2023stochasticinterpolants}
\def\WFkoehler2024Key{koehler2024flowmatching}
\let\WForigCiteTransRelax\WForiginalCite
\def\WFcitePlainTransRelax#1{%
  \def\WFciteTempKeyTransRelax{#1}%
  \ifx\WFciteTempKeyTransRelax\WFcuturi2013Key Cuturi, 2013%
  \else\ifx\WFciteTempKeyTransRelax\WFlipman2023Key Lipman et al., 2023%
  \else\ifx\WFciteTempKeyTransRelax\WFalbergo2023Key Albergo \& Vanden-Eijnden, 2023%
  \else\ifx\WFciteTempKeyTransRelax\WFkoehler2024Key Koehler, Heckett \& Risteski, 2024%
  \else \WForigCiteTransRelax{#1}%
  \fi\fi\fi\fi
}
\let\cite\WFcitePlainTransRelax
\makeatother

\subsection{Transport Relaxation for reward shaping}
\label{sec:method:transport-relaxation}

The reward-shaping layer of \S\ref{sec:end-to-end} aggregates seven
channels (qed, sa, binding, novelty, reinvent4, logp, ring\_count)
into a single scalar through the \texttt{RewardAggregator}; the
syntheisis-side and dual-system components of
\S\ref{sec:deflex} then re-shape the aggregated reward into a
\emph{learned shaping} term that closes the F5 channel of the
SynFlowNet trajectory-balance loss.  This subsection re-grounds that
shaping term in a \emph{transport-relaxation} objective:
Jiang-Xu-Zhang 2026 prove that a relaxed OT problem with a soft
endpoint penalty recovers the Schr\"odinger Bridge solution as the
penalty coefficient $\to 0$ and stays differentiable throughout
\cite{jiang2026sb_relax}.  The sub-section is the
theoretical-method reference for the W2.2 close-out that promotes
Theorem~2 of the R18 backlog from a paper-claim to a mathematical
skeleton + adapter pair.

\paragraph{Why a relaxed transport problem.}
A standard entropic OT problem (Cuturi 2013, Sinkhorn
\cite{cuturi2013sinkhorn}) enforces the marginal constraints
$P \mathbf{1} = a$, $P^{\top} \mathbf{1} = b$ \emph{exactly}; the
constraint is hard, and the trajectory-balance loss of the SynFlowNet
is not differentiable through the Sinkhorn fixed-point in the small-$
\varepsilon$ regime.  The relaxed version replaces the marginals by
soft blends:
\begin{equation}
  \label{eq:relaxed-marginals}
  \tilde{a} \;=\; (1 - \alpha)\,a + \alpha\, \mathbb{1}/n, \qquad
  \tilde{b} \;=\; (1 - \alpha)\,b + \alpha\, \mathbb{1}/m, \qquad
  \alpha \;=\; \frac{\lambda_{\mathrm{relax}}}{1 + \lambda_{\mathrm{relax}}},
\end{equation}
and solves the same entropic OT problem with these soft marginals.
Standard OT is recovered at $\lambda_{\mathrm{relax}} = 0$;
$\lambda_{\mathrm{relax}} \to \infty$ collapses the coupling to the
uniform plan $P_{ij} = 1/(nm)$.

\paragraph{Honest framing.}
The math layer is \textsc{shipped} as a CPU-only mathematical
skeleton at
\texttt{molmetal/math\_skills/transport\_relaxation.py} (25/25 unit
tests in
\texttt{molmetal/math\_skills/tests/test\_transport\_relaxation.py}).
The W2.2 adapter
(\texttt{molmetal/adapters/transport\_relaxation\_integration.py}) is
\textsc{shipped} and exposes
\texttt{augment\_learned\_shaping(learned\_shaping,
lambda\_relax)} as a drop-in wrapper for
\texttt{molmetal/molmetal\_lam/reward/learned\_shaping.py}.  The
projected pK$_{\mathrm{bind}}$ MAE lift ($0.42 \to 0.30$) is
\PROJECTED, \emph{not} \MEASURED; the $\lambda_{\mathrm{relax}}$
sensitivity sweep over $[10^{-4},\,10^{-2}]$ is the W2.2 follow-up
on the R18 Phase 2 backlog.

\subsubsection{Mathematical statement (informal)}
\label{sec:method:transport-relaxation-statement}

Let $\mu_0, \mu_1$ be two probability measures on a state space
$\mathcal{X}$, and let $\mathcal{P}(\mu_0, \mu_1)$ denote the set of
couplings.  The \emph{transport-relaxation problem} is
\begin{equation}
  \label{eq:transport-relaxation-problem}
  \pi^{\star}_{\alpha} \;=\; \arg\min_{\pi \in \mathcal{P}(\mu_0, \mu_1)}
        \int_{\mathcal{X} \times \mathcal{X}} c(x, y) \,\mathrm{d}\pi(x, y)
        \;+\; \varepsilon \, \mathrm{KL}\!\bigl(\pi \,\|\, \mu_0 \otimes \mu_1\bigr)
        \;+\; \lambda_{\mathrm{relax}}\, \mathrm{KL}\!\bigl(\pi_{T} \,\|\, q_{T}\bigr),
\end{equation}
where $\pi_{T}$ is the time-$T$ marginal of $\pi$, $q_{T}$ is a
reference Boltzmann distribution over docked poses (the
\texttt{BoltzmannEndpoint} of the math skeleton), and $\alpha =
\lambda_{\mathrm{relax}} / (1 + \lambda_{\mathrm{relax}})$.
Theorem~3.1 of \cite{jiang2026sb_relax} proves
$\pi^{\star}_{\alpha} \to \pi^{\star}_{\mathrm{SB}}$ as
$\lambda_{\mathrm{relax}} \to 0$; Theorem~3.2 proves the relaxed
problem stays differentiable in $\lambda_{\mathrm{relax}}$ throughout
the sweep.

\paragraph{Boltzmann endpoint.}
The reference $q_{T}$ is the canonical Boltzmann distribution over
docked poses:
\begin{equation}
  \label{eq:boltzmann-endpoint}
  q_{T}(y) \;\propto\; \exp\!\bigl(- \beta\, E_{\mathrm{Vina}}(y)\bigr),
\end{equation}
where $E_{\mathrm{Vina}}(y)$ is the QuickVina docking score at pose
$y$ and $\beta = 1/(k_{\mathrm{B}} T_{\mathrm{eff}})$ is the effective
inverse temperature.  $q_{T}$ is the \emph{right} reference for a
relaxed transport problem because it concentrates on the
low-energy docked poses; the transport-relaxation KL term
$\mathrm{KL}(\pi_{T} \,\|\, q_{T})$ then pushes the time-$T$
marginal of the coupling towards the docked-pose distribution.

\paragraph{Relationship to GFlowNet TB loss.}
The SynFlowNet trajectory-balance (TB) loss of
\S\ref{sec:deflex} decomposes as
$\mathcal{L}_{\mathrm{TB}} = \mathcal{L}_{\mathrm{TB}}^{\mathrm{base}}
+ \lambda_{\mathrm{relax}}\, \mathrm{KL}(p_{T} \,\|\, q_{T})$, where
$p_{T}$ is the marginal pose distribution at the leaf state and
$q_{T}$ is the Boltzmann endpoint.  Theorem~3.2 of
\cite{jiang2026sb_relax} guarantees the gradient
$\nabla_{\theta} \mathcal{L}_{\mathrm{TB}}$ is well-defined for any
$\lambda_{\mathrm{relax}} > 0$; the hard-marginal regime
($\lambda_{\mathrm{relax}} = 0$) is recovered as a limit.

\subsubsection{Mol-Metal surface: learned\_shaping.py augmentation}
\label{sec:method:transport-relaxation-surface}

The three primitives in
\texttt{molmetal/math\_skills/transport\_relaxation.py} are the
mathematical backbone that the W2.2 adapter wraps:

\begin{enumerate}
  \item \textbf{\texttt{RelaxedOT}} (lines ~280--620).  Computes the
        relaxed OT coupling of Eq.~\eqref{eq:transport-relaxation-problem}
        via log-domain Sinkhorn on the soft marginals
        $\tilde{a}, \tilde{b}$ of Eq.~\eqref{eq:relaxed-marginals}.

  \item \textbf{\texttt{BoltzmannEndpoint}} (lines ~640--820).  The
        canonical Boltzmann distribution over docked poses of
        Eq.~\eqref{eq:boltzmann-endpoint}, parameterised by the
        QuickVina energy $E_{\mathrm{Vina}}$ and the effective
        inverse temperature $\beta$.

  \item \textbf{\texttt{TransportRelaxationTerm}} (lines ~840--1280).
        The top-level entry point.  Computes
        $\lambda_{\mathrm{relax}}\, \mathrm{KL}(p_{T} \,\|\, q_{T})$
        and exposes a \texttt{sweep\_lambda} helper for the standard
        sensitivity-analysis sweep over
        $[10^{-4},\,10^{-2}]$.
\end{enumerate}

\paragraph{W2.2 adapter.}
The W2.2 adapter at
\texttt{molmetal/adapters/transport\_relaxation\_integration.py}
wraps the math layer with an
\texttt{augment\_learned\_shaping(learned\_shaping, lambda\_relax)}
function (line ~140) that adds
$\lambda_{\mathrm{relax}}\, \mathrm{KL}(p_{T} \,\|\, q_{T})$ to the
existing \texttt{learned\_shaping.py} reward term.  The augmentation
is a drop-in wrapper: it does not modify the live
\texttt{learned\_shaping.py} module (the R17 close-out is locked)
and returns a new shaping term with the transport-relaxation KL
appended.

\paragraph{Per-state aggregation.}
The reward augmentation is applied at the MCTS leaf-state level: for
each closed $\beta$-NF candidate (a typed term satisfying the
$\Phi$-refinement of \S\ref{sec:metal-geometry-prior}), the
\texttt{RewardAggregator} computes the seven-channel scalar
$\mathcal{R}_{\mathrm{7ch}}$ and the
\texttt{augment\_learned\_shaping} wrapper appends the transport-
relaxation KL.  The final SynFlowNet advantage is
$\mathcal{A} = \mathcal{R}_{\mathrm{7ch}} - \lambda_{\mathrm{relax}}\,
\mathrm{KL}(p_{T} \,\|\, q_{T})$.

\subsubsection{Expected gain and honest framing}
\label{sec:method:transport-relaxation-expected}

\paragraph{Projected metric gain.}
The R18 backlog projects a pK$_{\mathrm{bind}}$ MAE lift of $0.42 \to
0.30$ (RDKit baseline $1.10$, \textsc{P0-2} of
\texttt{mol-metal-r18-math/SKILL.md} \S\,priority map).  The gain
comes from three sources:

\begin{itemize}
  \item the Boltzmann endpoint $q_{T}$ concentrates probability mass
        on the low-energy docked poses, providing a sharper
        reference than the unconstrained uniform marginal;
  \item the relaxed marginals keep the Sinkhorn fixed-point
        well-conditioned throughout the $\lambda_{\mathrm{relax}}$
        sweep, avoiding the numerical blow-up of the hard-marginal
        regime;
  \item the soft-endpoint penalty is differentiable in
        $\lambda_{\mathrm{relax}}$, so the GFlowNet TB loss back-
        propagates through the entire chain end-to-end.
\end{itemize}

\paragraph{Honest framing on the pK$_{\mathrm{bind}}$ MAE lift.}
The $-0.12$ MAE reduction is \PROJECTED, not \MEASURED.  The live
reward pipeline has the \texttt{augment\_learned\_shaping} hook in
place (the W2.2 close-out ships the adapter) but the
$\lambda_{\mathrm{relax}}$ sensitivity sweep is on the R18 Phase 2
backlog.  The math + adapter layers are CPU-only and \textsc{shipped};
the projected lift is grounded in the Theorem 2 chain but is not yet
exercised on the live CrossDocked2020 pK$_{\mathrm{bind}}$
prediction task.

\paragraph{Frontier anchors.}
The transport-relaxation chain sits on top of the flow-matching
core anchors \cite{lipman2023flowmatching,albergo2023stochasticinterpolants,
koehler2024flowmatching} and the original Sinkhorn divergence
\cite{cuturi2013sinkhorn}; the 24-paper Frontier survey is
consolidated in
\texttt{molmetal/reports/wf\_cfm\_m3\_m5/frontier\_24\_papers.md}.

\paragraph{Cross-section pointers.}
The transport-relaxation term augments the \texttt{learned\_shaping}
F5 channel of \S\ref{sec:deflex} (the dual-system Deflex reward
pipeline); the Boltzmann endpoint $q_{T}$ reuses the QuickVina
score from the parallel validation band (stages (e) of
Figure~\ref{fig:pipeline}); the $\lambda_{\mathrm{relax}}$ sensitivity
sweep is documented in the W2.2 close-out at
\texttt{molmetal/reports/wf\_transport\_relaxation\_w22/final.md}.
The §4 round-17 honest-negative framing retains the 0/192 decode
ratio baseline; the transport-relaxation term is the Phase 2 upgrade
target for the reward-shaping layer.

% -----------------------------------------------------------------------------
% Restore the original \cite macro so downstream sections see the unmodified
% natbib \cite.  Same pattern as 03_3_pll_heuristic.tex:295-297.
% -----------------------------------------------------------------------------
\makeatletter
\let\cite\WForigCiteTransRelax
\makeatother